\subsection{伪度量}

定义 1. 若 X 是一个集合，X 上的伪度量是指 X × X 到扩充实直线 R 的区间 {[}0, +∞{]} 中的、满足下列条件的任一映射 f：

\((\mathrm{EC}_{\mathbf{I}})\) 对一切 \(x \in \mathbf{X}, f(x, x) = 0\) 。

\((\mathrm{EC}_{\mathrm{II}})\) 对一切 \(x \in \mathbf{X}\) 与一切 \(y \in \mathbf{X}, f(x, y) = f(y, x)\) （对称性）。\((\mathrm{EC}_{\mathrm{III}})\) 对 \(x, y, z\) —— \(\mathbf{X}\) 中的任意元素 —— ，

\[
f (x, y) \leqslant f (x, z) + f (z, y)
\]

（三角不等式）。

例。1) 在实数空间 \(\mathbf{R}^n\) 上，欧氏距离（第六章，§ 2，第 1 号）是一个伪度量。

\begin{enumerate}
\def\labelenumi{\arabic{enumi})}
\setcounter{enumi}{1}
\item
  若 X 是任一集合，则在 \(X \times X\) 上由条件“\(f(x, x) = 0\) 对一切 \(x \in X\) 成立，而 \(f(x, y) = +\infty\) 当 \(x \neq y\) 时”所定义的函数 f 是 X 上的一个伪度量。
\item
  若 \(\mathbf{X}\) 是任一集合且 \(g\) 是定义在 \(\mathbf{X}\) 上的任一有限实值函数，则函数 \(f\) —— 在 \(\mathbf{X} \times \mathbf{X}\) 上由 \(f(x, y) = |g(x) - g(y)|\) 定义的 —— 是 \(\mathbf{X}\) 上的一个伪度量。
\end{enumerate}

* 4) 设 X 是 R 的区间 {[}0, 1{]} 到 R 中的所有连续映射所成的集。若对 X 的每对元素 x, y 令

\[
f (x, y) = \int_ {0} ^ {1} | x (t) - y (t) | d t,
\]

则 \(f\) 是 \(\mathbf{X}\) 上的一个伪度量。*

注。1) 上面的例 2 表明，伪度量对 X 的某些元素对可以取值 +∞。

\begin{enumerate}
\def\labelenumi{\arabic{enumi})}
\setcounter{enumi}{1}
\tightlist
\item
  若 \(f\) 是 \(\mathbf{X}\) 上的一个伪度量，则一般说来可能有 \(f(x, y) = 0\) 而 \(x = y\) 不成立，如上面的例 3 所示（参看§ 2）。
\end{enumerate}

由三角不等式可知，若 \(f(x, z)\) 与 \(f(y, z)\) 有限，则 \(f(x, y)\) 也有限；而且此时有

\[
f (x, z) \leqslant f (y, z) + f (x, y) \quad \text { and } \quad f (y, z) \leqslant f (x, z) + f (x, y),
\]

从而

\[
\left| f (x, z) - f (y, z) \right| \leqslant f (x, y).\tag{1}
\]

若 \(f\) 是 \(\mathbf{X}\) 上的一个伪度量，则 \(\lambda f\) （对任何有限实数 \(\lambda > 0\) ）也是伪度量。若 \((f_t)_{t \in I}\) 是 \(\mathbf{X}\) 上伪度量的任一族，则和 \(\sum_{i \in I} f_i(x, y)\) 对一切 \((x, y) \in \mathbf{X} \times \mathbf{X}\) 有定义；若用 \(f(x, y)\) 表示这个和的值，则 \(f\) 是 \(\mathbf{X}\) 上的一个伪度量。同样，上包络 \(g\) —— 族 \((f_t)\) 的上包络 —— （第四章，§ 5，第 5 号）是 \(\mathbf{X}\) 上的一个伪度量，因为关系 \(f_i(x, y) \leqslant f_i(x, z) + f_i(z, y)\) 蕴含

\[
\sup _ {t \in I} f _ {t} (x, y) \leqslant \sup _ {t \in I} \left(f _ {t} (x, z) + f _ {t} (z, y)\right) \leqslant \sup _ {t \in I} f _ {t} (x, z) + \sup _ {t \in I} f _ {t} (z, y)
\]

{[}Chapter IV, § 5, no. 7, formula (17){]}.

